% Chapter 9 reference resources. Included by ch09_resources.tex.
\section{Reference Resources}
\label{sec:comp-ch09-reference}

This reference section distinguishes concepts, methods, reporting guidance,
standards, regulations, and author-developed tools used with Chapter~9. The
categories should not be treated as interchangeable. A law establishes obligations
within its scope; a standard or framework structures practice; a reporting guideline
supports transparency; an appraisal tool supports critical review; and an
institutional template records a local decision. None establishes clinical or
economic value by name alone. 

\subsection{Key Terms}
\begin{longtable}{@{}p{0.27\textwidth}p{0.65\textwidth}@{}}
\caption{Key terms for Chapter 9}
\label{tab:comp-ch09-key-terms}\\
\toprule
\textbf{Term} & \textbf{Working definition and decision relevance} \\
\midrule
\endfirsthead
\toprule
\textbf{Term} & \textbf{Working definition and decision relevance} \\
\midrule
\endhead
\bottomrule
\endlastfoot
Problem-first procurement & A process that defines need, comparator, evidence, and safeguards before product selection. \\
Mandatory safeguard & A non-compensatory requirement that must be met before weighted comparison. \\
Vendor governance & Continuing oversight of evidence, performance, change, continuity, and exit after award. \\
\end{longtable}

\subsection{Status and Appropriate Use of Selected Instruments}
\begin{longtable}{@{}p{0.25\textwidth}p{0.23\textwidth}p{0.44\textwidth}@{}}
\caption{Selected instruments relevant to Chapter 9}
\label{tab:comp-ch09-instrument-status}\\
\toprule
\textbf{Instrument} & \textbf{Status or type} & \textbf{Appropriate use and limitation} \\
\midrule
\endfirsthead
\toprule
\textbf{Instrument} & \textbf{Status or type} & \textbf{Appropriate use and limitation} \\
\midrule
\endhead
\bottomrule
\endlastfoot
NICE Evidence Standards Framework & Evidence guidance & Supports evidence planning for digital health technologies. \\
DTAC & Jurisdiction-specific assessment criteria & Supports baseline digital-technology due diligence in England. \\
EU model contractual clauses & Public procurement resource & Provides adaptable AI contractual clauses. \\
FAIR-AI & Implementation and review framework & Links pre-implementation review with post-implementation monitoring. \\
\end{longtable}

\subsection{Method and Evidence Selection}
Select resources according to the question being answered:
\begin{itemize}
 \item use a reporting guideline to make a study transparent, not to certify that the design was appropriate;
 \item use an appraisal tool to identify risk of bias and applicability concerns, not to replace local judgement;
 \item use an implementation framework for a defined determinant, process, outcome, or context question;
 \item use a standard or management system to structure repeatable responsibilities, while verifying the current edition and legal context;
 \item use author-developed tools as proposed decision aids, with explicit status, version, limitations, and validation requirements; and
 \item retain a versioned record of why each resource was selected, how it was adapted, and which decision it informed.
\end{itemize}

\subsection{Formula and Interpretation Note}
The following relationship is used in Chapter~9 or its companion resources:
\begin{equation}
S_v=\sum_{d=1}^{D}w_ds_{vd},\qquad G_v=\prod_{k=1}^{K}g_{vk}
\label{eq:comp-ch09-reference-formula}
\end{equation}
Its variables, perspective, time horizon, uncertainty, and limitations must be
defined before use. A formula should not be added merely for symmetry with other
chapters; it is useful only when it clarifies a decision and can be populated with
defensible evidence.

\subsection{Update Discipline}
Before operational use, verify the current edition, jurisdictional status, scope,
and evidence base of every external instrument. Record the access date for online
policy material, preserve the exact version used, and reassess applicability after a
material change in law, regulation, technology, intended use, evidence, or service
context.
